% Generated draft: requires mathematical and source review.
\documentclass[11pt]{article}
\usepackage[margin=25mm]{geometry}
\usepackage{fontspec}
\setmainfont{Latin Modern Roman}
\usepackage{amsmath,amssymb,mathtools,unicode-math}
\setmathfont{Latin Modern Math}
\usepackage{enumitem,needspace,xurl,hyperref}
\setlength{\parindent}{0pt}
\setlength{\parskip}{5pt}
\setlength{\emergencystretch}{3em}
\newcommand{\N}{\mathbb{N}}
\newcommand{\Z}{\mathbb{Z}}
\newcommand{\Q}{\mathbb{Q}}
\newcommand{\R}{\mathbb{R}}
\newcommand{\C}{\mathbb{C}}
\newcommand{\F}{\mathbb{F}}
\newcommand{\A}{\mathbb{A}}
\newcommand{\PP}{\mathbb{P}}
\newcommand{\E}{\mathbb{E}}
\newcommand{\eps}{\varepsilon}
\newcommand{\dd}{\,\mathrm{d}}
\begin{document}
\section*{UnsolvedMath ID: 30004951}
\textbf{Source problem code:} catalogue entry
\subsection{Definitions and mathematical statement}
Let \(n\geq 2\) be an integer and let \(k\in\{1,\ldots,n-1\}\). For a polytope \(P\subset\R^n\), let \(V_n(P)\) denote its \(n\)-dimensional Lebesgue volume. When \(V_n(P)>0\), define its volume centroid by
\[
g(P)=\frac{1}{V_n(P)}\int_P x\,\dd x.
\]
The supplied record refers to volume-decomposition functionals \(X_k\), but supplies neither their definitions nor the decomposition identity mentioned in its original excerpt. Consequently, \(X_k(P)\) below is an unresolved source-specific quantity, whose definition and normalization must be recovered before this becomes a self-contained mathematical statement.

The question in the supplied clean statement is whether
\[
\forall n\in\Z,\ n\geq2,\quad
\forall k\in\{1,\ldots,n-1\},\quad
\exists c(n,k)\in\R\quad
\forall P\subset\R^n\text{ a polytope with }V_n(P)>0\text{ and }g(P)=0,
\qquad
\frac{X_k(P)}{V_n(P)}\leq c(n,k).
\]
Here the constant must be finite and depend only on \(n\) and \(k\). The requirement \(V_n(P)>0\), equivalently full dimensionality for a polytope in \(\R^n\), is necessary for this quotient and the displayed centroid formula; it is not stated explicitly in the supplied problem and requires confirmation against the intended source conventions.

\textit{Unresolved normalization.} The supplied original excerpt instead asks for \(X_k(P)\leq c(n,k)\), without a volume denominator. No definition or normalization is supplied that reconciles these formulations. The displayed quantified question records the clean statement's target, without asserting that the original excerpt is equivalent to it.

\textit{Supplied context.} The excerpt relates the decomposition functionals to the Lutwak--Yang--Zhang \(U\) functional and describes a 2001 conjecture concerning \(X_n\) and volume, with equality precisely for parallelotopes, as subsequently settled. Its displayed inequality is corrupted in the extraction and cannot be restated reliably from this record. The excerpt also says that Problem \(X\) in \(\R^3\) was satisfactorily solved by Liu--Sun--Xiong; the scope of that assertion relative to the normalized question above is not specified.
\subsection{Short English statement}
For each dimension \(n\geq2\) and each \(1\leq k<n\), the clean statement asks whether \(X_k(P)/V_n(P)\) has a finite upper bound depending only on \(n,k\) over polytopes with volume centroid at the origin. The definition of \(X_k\) is missing, and the original excerpt omits the volume denominator. Positive volume is required to interpret the quotient. These issues require human mathematical review.
\subsection{Sources}
\begin{itemize}
\item Supplied local UnsolvedMath registry record, permanent ID \(30004951\), exact source problem code \texttt{catalogue entry}, title ``Uniform Bounds for Volume-Decomposition Functionals.'' The record belongs to the Oberwolfach Reports Open Problems collection and gives proposed year 2021; no individual proposer is supplied.
\item Supplied report citation: \emph{Convex Geometry and its Applications} (2022), Oberwolfach Reports, pp.\ 3187--3266, DOI 10.4171/owr/2021/59. \url{https://doi.org/10.4171/owr/2021/59}. The background describes these reports as documenting mathematical workshops and their open-problem discussions. This reference is supplied provenance and has not been independently consulted for this draft.
\item Supplied literature source: ``Liu--Sun--Xiong on X sub 2 and X sub 3,'' \url{https://arxiv.org/abs/2404.09152}. No further bibliographic details are supplied.
\item Supplied literature source: ``original OWR report,'' \url{https://doi.org/10.4171/owr/2021/59}.
\item Literature assessment dated exactly \texttt{2026-08-22}: the supplied curation labels the problem open, says the volume-decomposition question is solved in dimension three and sharp inequalities are known for \(X_2\) and \(X_3\), and says the extracted all-\(k\) boundedness question is not settled in general. This is supplied dated curation, not an independently verified literature assessment. The background explicitly requests independent verification of time-sensitive status and cited claims.
\item The registry labels the statement \texttt{corrected\_verified} and reports that it was checked against the official Oberwolfach report and preserves its hypotheses, notation, and open target. These are supplied verification claims; the missing definitions and normalization discrepancy remain unresolved in the evidence provided here. The original excerpt's numbered references [1]--[5], including its Liu--Sun--Xiong reference [3], lack a supplied bibliography and cannot be identified reliably from their numbers alone.
\end{itemize}
\end{document}
